% META: {"id":"1SPE-VARALEA-EX-033","chapitre":"1SPE-VARIABLES-ALEATOIRES","type_objet":"exercice","capacites_codes":[],"extension_codes":["X2"],"programme_alignment":"OPTIONAL_EXTENSION","extension_label":"Approfondissement — Vers la Terminale","parcours":2,"competences":["calculer","raisonner"],"duree_min":15,"mode_creation":"ex_nihilo","sources_inspiration":[],"fichier_tex":"chapitres/1SPE-VARIABLES-ALEATOIRES/exercices/1SPE-VARALEA-EX-033.tex","corrige_tex":"chapitres/1SPE-VARIABLES-ALEATOIRES/corriges/1SPE-VARALEA-CO-033.tex","status":"approved","origin":"generated","status_history":["generated","audited","accepted","approved"],"release_acceptance":"RELEASE_OWNER_BATCH_ACCEPTANCE","acceptance_closure_digest":"sha256:9b3ccf9a81c5520fbb7e03b7b2d3e7bf2057a02834b6d908bf3be5f49f3b553f","acceptance_receipt_digest":"sha256:4fad86f0282e0fa4e0419cca1ad6379ae7af5a280c04a4a90880f90a1c044242"}
\begin{center}\textbf{Approfondissement — Vers la Terminale}\end{center}
% BEGIN-VERIFY
% from sympy import Rational
% # Retrouver n et p a partir de E=6 et V=4.2
% E = 6; V = Rational(42,10)
% p = 1 - V/E
% assert p == Rational(3,10)
% n = E/p
% assert n == 20
% END-VERIFY
\begin{exercice}{1SPE-VARALEA-EX-033}{2}{15}
On sait que $X \sim \mathcal{B}(n, p)$ avec $E(X) = 6$ et $V(X) = 4{,}2$. Déterminer $n$ et $p$.
\end{exercice}
